\section{相标记}
\label{sec:markers}

\subsection{结构因子与 AFM 标记}
\label{sec:sq}

为识别反铁磁相，我们使用结构因子
\begin{equation}
S(q)=\frac{1}{L}\sum_{i,j}e^{iq(i-j)}\langle Z_iZ_j\rangle,
\end{equation}
其中 $\langle Z_iZ_j\rangle$ 为基态关联子，本节恒取 $L=8$。
即 $S(q)$ 是 N\'eel 型关联的傅里叶权重，长程反铁磁序表现为 $q=\pi$ 处的
Bragg 型尖峰。三相各一个代表点处的 $S(q)$ 曲线（相标签与坐标见数据集），
共用 $q$ 网格（$99$ 点，含 $q=\pi$），见图~\ref{fig:D03}。
在面向真机的测量协议中，使用 $q=\pi$ 分量作为 AFM 诊断量；
其对不同相的实际区分结果放在结果部分报告，而不在方法部分预先断言。

\begin{figure}[htbp]
\centering
\includegraphics[width=\columnwidth]{exp02_D03.pdf}
\caption{共用 $q$ 网格上三相各一个代表点处的结构因子 $S(q)$；
只有 AFM 点在 $q=\pi$ 处见峰。
\label{fig:D03}}
\end{figure}
